% Figure from MATH 556 notes, section 1. Compile with the notes preamble.
\begin{tikzpicture}[scale=1.25, >=stealth]
    % the classes x + Y: lines parallel to Y (slope 0.35), offsets 0.9k
    \foreach \k in {-1,1,2} {
        \draw[figR!70, thick] (-2.2,{-0.77+0.9*\k}) -- (2.6,{0.91+0.9*\k});
    }
    \node[figR, anchor=west, font=\scriptsize] at (2.6,1.81) {$[x]$};
    \node[figR, anchor=west, font=\scriptsize] at (2.6,2.71) {$[2x]$};
    \node[figR, anchor=west, font=\scriptsize] at (2.6,0.01) {$[-x]$};
    % Y itself = the zero class
    \draw[very thick] (-2.2,-0.77) -- (2.6,0.91) node[anchor=west, font=\scriptsize] {$Y = [0]$};
    % complement W (blue), direction (0.35,1)
    \draw[figB, thick] ({-1.3*0.35},-1.3) -- ({2.75*0.35},2.75) node[above, font=\scriptsize] {$W$};
    % second complement W' (green), direction (-0.6,1)
    \draw[figG, thick] ({1.6*0.6},-1.6) -- ({-2.6*0.6},2.6) node[above, font=\scriptsize] {$W'$};
    % intersection points with the classes
    \foreach \k in {-1,1,2} {
        \fill[figB] ({1.0256*\k*0.35},{1.0256*\k}) circle (1.3pt);
        \fill[figG] ({-0.7438*\k*0.6},{0.7438*\k}) circle (1.3pt);
    }
    \node[figB, anchor=west, font=\scriptsize] at (0.40,0.93) {$w$};
    \node[figG, anchor=south, font=\scriptsize, inner sep=1pt] at (-0.35,0.88) {$w'$};
    \fill (0,0) circle (1.1pt);
\end{tikzpicture}
